\documentclass[12pt,a4paper]{article}

% 数学核心宏包（ctexart已内置基础数学，这里补充高级功能）
\usepackage{amsmath, amssymb, amsthm, mathtools}

% 页面与列表
\usepackage{geometry}
\geometry{margin=1in}
\usepackage{enumitem}

% 常用集合符号快捷命令
\newcommand{\R}{\mathbb{R}}
\newcommand{\Z}{\mathbb{Z}}
\newcommand{\N}{\mathbb{N}}
\newcommand{\Q}{\mathbb{Q}}
\newcommand{\cart}{\times}
\newcommand{\inter}{\cap}
\newcommand{\union}{\cup}

% 定理/命题环境
\newtheorem*{proposition}{Proposition}
\newtheorem*{example}{Example}
\newtheorem*{definition}{Definition}

\begin{document}
\begin{center}
\LARGE \textbf{MTH024 Quiz1}
\end{center}

\noindent{\Large \textbf{Problem 1}}\\
$A \cup B = \{x|x \in A \;\ or \;\ x \in B\} = \{1,2,3,4,5,6,7\}$\\
$A \cap B = \{x|x \in A \;\ and \;\ x \in B\}$ =\{4,5\}

\vspace{1.5em}

\noindent{\Large \textbf{Problem 2}}\\
$Define:E^c $ is set of odd nimbers\\
$E^c = \{x\in \mathbb{N} | x\;\ is\;\ odd\;\ numbers \}=\{x\in \mathbb{N} | x=2k+1,k \in \mathbb{N} \cap \{0\}\}$

\vspace{1.5em}
\noindent{\Large \textbf{Problem 3}}\\
$ x \in A \cap(B \cup C  ) \Leftrightarrow x \in A \;\ and \;\ x \in B \cup C $\\
$\Leftrightarrow x \in A \;\ and \;\ (x\in B \;\ or \;\ x \in C) $\\
$\Leftrightarrow (x \in A \;\ and \;\ x \in B) \;\ or \;\ (x \in A \;\ and \;\ x \in C)$\\
$\Leftrightarrow x \in A \cap B \;\ or \;\ x \in A \cap C $\\
$\Leftrightarrow x \in (A \cap B)\cup (A \cap C)$

\vspace{1.5em}
\noindent{\Large \textbf{Problem 4}}\\
$(a):\mathbb{Z} \;\ \mathbb{R} \;\ \mathbb{Q} $\\
$(b):\mathbb{Q} \;\ \mathbb{R} $\\
$(c):\mathbb{R} $\\
$(d):\mathbb{Q} \;\ \mathbb{R} $

\vspace{1.5em}
\noindent{\Large \textbf{Problem 5}}\\
let $a$ be an even number, $a=2k$, $b$ is any integer\\
$a\times b=2k\times b=2(kb)$\\
$\because 2(kb) \;\ is \;\ even$,\\
$\therefore a\times b$ is even.

\vspace{1.5em}
\noindent{\Large \textbf{Problem 6}}\\
Suppose for contradiction that $n$ is an odd number, such that $n^2$ is even.\\
$\because n=2k+1$\\
$\therefore n^2=4k^2+4k+1=2(2k^2+2k)+1$\\
$\because 2(2k^2+2k)$ is an even number\\
$\therefore 2(2k^2+2k)+1$ is an odd number \\
But this is not true since $n^2$ is even.\\
Hence if $n^2$ is an even integer, then $n$ must also be an even integer.

\vspace{1.5em}
\noindent{\Large \textbf{Problem 7}}\\
Suppose for contradiction that $ \sqrt{2} $ is an rational number, \\
Exist coprime integers $p$ and $q(q\neq 0)$ such that $ \sqrt{2}=\frac{p}{q} $.\\
$\therefore 2=\frac{p^2}{q^2} \Rightarrow p^2=2q^2$\\
$\therefore p$ is even, let $p=2k$\\
$\because (2k)^2=2q^2\Leftrightarrow q^2=2k^2$\\
$\therefore q$ is even.\\
Since $p$ and $q$ are both even numbers, $p$ and $q$ are not coprime. 
But it is not true since $p$ and $q$ should be coprimed, then $ \sqrt{2} $ is an irrational number.

\vspace{1.5em}
\noindent{\Large \textbf{Problem 8}}\\
base: when $n=1$, $1=\frac{1\times (1+1)}{2}$ satisfies the function.\\
induction: when $n=k$, $1+2+\dots +k=\frac{k(k+1)}{2}$,\\
At the same time, when $n=k+1$, $1+2+\dots +k+1=\frac{k(k+1)}{2}+(k+1)$\\
$\because \frac{k(k+1)}{2}+(k+1)= \frac{(k+1)(k+2)}{2}$ \\
$\therefore 1+2+\dots +n=\frac{n(n+1)}{2}$ is proved.

\vspace{1.5em}
\noindent{\Large \textbf{Problem 9}}\\
base: when $n=1$, $1^2=\frac{1\times (1+1)\times (2+1)}{6}$ satisfies the function.\\
induction: when $n=k$, $1^2+2^2+\dots +k^2=\frac{k(k+1)(2k+1)}{6}$,\\
At the same time, when $n=k+1$, $1^2+2^2+\dots +(k+1)^2=\frac{k(k+1)(2k+1)}{6}+(k+1)^2$\\
$\because \frac{k(k+1)(2k+1)}{6}+(k+1)^2=\frac{(k+1)(k+2)(2k+3)}{6}=\frac{(k+1)(k+2)(2(k+1)+1)}{6}$\\
$\therefore 1^2+2^2+\dots +n^2=\frac{k(k+1)(2k+1)}{6}$ is proved.

\vspace{1.5em}
\noindent{\Large \textbf{Problem 10}}\\
base: when $n=4$, $4!=24>2^4=16$ satisfies $n^2>2^n$.\\
induction: when $n=k$, $k!>2^k$,\\
At this time, when $n=k+1$, $(k+1)!=(k+1)k!>(k+1)2^k$ \\
$\because n\geqslant 4$\\
$\therefore k+1\geqslant  5 >2$\\
$\because (k+1)2^k>2\times 2^k=2^{k+1}$\\
$\therefore (k+1!)>2^{k+1}$\\
$\therefore n!>2^n$ is proved.
\end{document}